\documentclass[conference]{IEEEtran}

\usepackage{iftex}
\ifPDFTeX
\usepackage[utf8]{inputenc}
\usepackage[T1]{fontenc}
\fi
\usepackage{amsmath,amssymb,amsthm,bm}
\usepackage{booktabs}
\usepackage{mathtools}
\usepackage{siunitx}
\usepackage[final]{microtype}
\usepackage[hidelinks]{hyperref}

\newtheorem{theorem}{Theorem}
\newtheorem{corollary}{Corollary}
\newtheorem{remark}{Remark}
\newcommand{\vct}[1]{\bm{#1}}
\newcommand{\mat}[1]{\bm{#1}}
\newcommand{\R}{\mathbb{R}}
\newcommand{\Id}{\mat I}

\title{Physical-MSE Scaling of Dual Pseudo-Measurement Regularization\\
in an OU-Driven Marine Inertial MEKF}

\author{\IEEEauthorblockN{Mikhail S. Grushinskiy}
\IEEEauthorblockA{Independent Researcher, USA, Fair Lawn\\
Bareboat Necessities GitHub Project\\
Email: mgrouch@users.github.com}}

\begin{document}
\maketitle

\begin{abstract}
A marine inertial multiplicative extended Kalman filter (MEKF) with a latent
Ornstein--Uhlenbeck (OU) acceleration state uses zero pseudo-measurements on
both displacement and velocity to suppress low-frequency integration drift.
The implemented adaptation schedules the corresponding pseudo-measurement
standard deviations as $r_{p,0}\propto\sigma_{aw}\tau^2$ and
$r_{v,0}\propto\sigma_{aw}\tau$, followed by a common cadence normalization.
This note derives a reduced physical mean-square-error (MSE) model for the two
pseudo channels jointly rather than treating them as independent penalties.
The associated two-state Kalman--Bucy algebraic Riccati equation admits a
closed-form solution parameterized by a position regularization frequency and
a dimensionless velocity-to-position strength ratio.  A weak-regularization
expansion separates residual-acceleration drift from distortion of the physical
wave.  For a fixed-shape sea and an approximately sea-independent low-frequency
posterior acceleration-error intensity, joint MSE minimization predicts the
base laws
$r_{p,0}^\star\propto\sigma_a^{4/5}\tau^{12/5}$ and
$r_{v,0}^\star\propto\sigma_a^{4/5}\tau^{7/5}$.
The two laws differ from the implemented powers by the same weak common factor,
and their ratio satisfies
$(r_{p,0}^\star/r_{v,0}^\star)^2=(3/2)M_{-2}/M_0$.
Consequently, the theory predicts $r_{p,0}^\star/r_{v,0}^\star\propto\tau$,
which matches the implemented relative period scaling exactly.  Evaluation on
the finite-band reference spectra places the theoretical coefficient ratio near
the implemented value.  The derivation supplies testable adaptation exponents
without fitting them to the existing schedule, while its reduced displacement
objective does not capture the full attitude and bias coupling of the MEKF.
\end{abstract}

\begin{IEEEkeywords}
Kalman filter, inertial navigation, marine motion estimation,
Ornstein--Uhlenbeck process, pseudo-measurement, regularization, mean-square
error, heave estimation.
\end{IEEEkeywords}

\section{Introduction}

Low-frequency drift is the central difficulty in displacement estimation from
inertial acceleration.  The OU--II estimator considered here regularizes this
integration chain with two nonphysical zero pseudo-measurements: one on
world-frame displacement $p$ and one on world-frame velocity $v$.  The two
channels are applied on the same periodic pseudo-update tick.  Their configured
standard deviations are therefore not independent observation uncertainties;
they are design parameters that trade drift suppression against distortion of
real wave motion.

The deployed adaptation was selected empirically and has the base form
\begin{equation}
 r_{p,0}^{\mathrm{base}}=c_p\sigma_{aw}\tau^2,
 \qquad
 r_{v,0}^{\mathrm{base}}=c_v\sigma_{aw}\tau,
 \label{eq:implemented-base}
\end{equation}
with current coefficients $c_p=0.65$ and $c_v=1.30$.  Here $\tau$ is the OU
correlation time and $\sigma_{aw}$ is the stationary standard deviation of the
latent OU acceleration prior.  Both pseudo channels use the same
cadence-normalization factor
\begin{equation}
 r_{j}=r_{j,0}^{\mathrm{base}}
 \sqrt{\frac{T_0}{T_S}},
 \qquad j\in\{p,v\},
 \label{eq:cadence-normalization}
\end{equation}
where the pseudo-update period satisfies $T_S\propto\tau$ away from safety
clamps.  This normalization preserves the continuous-equivalent information
density associated with a given base value.

Equation~\eqref{eq:implemented-base} has the correct dimensions and performs
well empirically, but dimensional consistency alone does not identify the
powers of $\sigma_{aw}$ and $\tau$.  The purpose of this note is to derive those
powers from a physical displacement-MSE criterion.  The analysis follows the
same distinction used in the OU--III regularization study: the internal Kalman
covariance of a filter that treats a zero pseudo-measurement as if it were a
physical observation is not itself the physical wave-estimation error.  A
useful design objective must include both residual stochastic drift and the
bias introduced by suppressing the real wave signal.

\section{Implemented OU--II Regularization Structure}
\label{sec:implementation}

With gyroscope and accelerometer biases enabled, the complete OU--II MEKF has
local error coordinates containing attitude, gyroscope bias, velocity,
displacement, latent OU acceleration, and accelerometer bias.  The present
analysis isolates the low-frequency translational subproblem
\begin{equation}
 \dot p=v,\qquad \dot v=n_a,
 \label{eq:double-integrator}
\end{equation}
where $n_a$ denotes the residual acceleration error presented to the integration
chain after acceleration estimation.  At frequencies below its posterior
correlation rate, approximate $n_a$ as white with two-sided intensity
$q_{\mathrm{eff}}$, abbreviated below as $q$.

The position and velocity pseudo-measurements are both zero.  If their discrete
standard deviations are $r_p$ and $r_v$ and the common update period is $T_S$,
the corresponding continuous-equivalent measurement-noise densities are
\begin{equation}
 \rho_p=r_p^2T_S,
 \qquad
 \rho_v=r_v^2T_S.
 \label{eq:continuous-densities}
\end{equation}
Under Eq.~\eqref{eq:cadence-normalization},
\begin{equation}
 \rho_p=(r_{p,0}^{\mathrm{base}})^2T_0,
 \qquad
 \rho_v=(r_{v,0}^{\mathrm{base}})^2T_0,
 \label{eq:base-density-equivalence}
\end{equation}
so the theoretical scaling of the continuous densities maps directly to the
base adaptation laws.

\section{Exact Dual-Channel Riccati Reduction}
\label{sec:care}

Consider the continuous model
\begin{equation}
 \dot{\vct x}=\mat A\vct x+\mat B n_a,
 \qquad
 \vct x=\begin{bmatrix}p&v\end{bmatrix}^{\top},
 \label{eq:state-model}
\end{equation}
with
\begin{equation}
 \mat A=\begin{bmatrix}0&1\\0&0\end{bmatrix},
 \qquad
 \mat B=\begin{bmatrix}0\\1\end{bmatrix},
 \qquad
 \mat Q=q\mat B\mat B^{\top}.
\end{equation}
The two continuous pseudo-observations have $\mat C=\Id$ and
\begin{equation}
 \mat R_c=\operatorname{diag}(\rho_p,\rho_v).
\end{equation}
The stabilizing covariance satisfies the Kalman--Bucy algebraic Riccati
equation \cite{KalmanBucy1961,Jazwinski1970}
\begin{equation}
 \mat A\mat P+\mat P\mat A^{\top}+\mat Q
 -\mat P\mat R_c^{-1}\mat P=0.
 \label{eq:care}
\end{equation}

Define the two natural regularization frequencies
\begin{equation}
 \omega_p:=\left(\frac{q}{\rho_p}\right)^{1/4},
 \qquad
 \omega_v:=\left(\frac{q}{\rho_v}\right)^{1/2},
 \label{eq:natural-frequencies}
\end{equation}
and the dimensionless relative velocity strength
\begin{equation}
 \chi:=\left(\frac{\omega_v}{\omega_p}\right)^2.
 \label{eq:chi}
\end{equation}

\begin{theorem}[Dual-channel CARE solution]
For $q,\rho_p,\rho_v>0$, the stabilizing solution of
Eq.~\eqref{eq:care} is
\begin{equation}
 \boxed{
 \mat P=\frac{q}{1+\chi}
 \begin{bmatrix}
 \dfrac{\sqrt{2+\chi}}{\omega_p^3} & \dfrac{1}{\omega_p^2}\\[2.5mm]
 \dfrac{1}{\omega_p^2} & \dfrac{\sqrt{2+\chi}}{\omega_p}
 \end{bmatrix}.}
 \label{eq:P-solution}
\end{equation}
The corresponding continuous pseudo-measurement gain
$\mat K=\mat P\mat R_c^{-1}$ is
\begin{equation}
 \boxed{
 \mat K=
 \begin{bmatrix}
 \dfrac{\omega_p\sqrt{2+\chi}}{1+\chi} &
 \dfrac{\chi}{1+\chi}\\[2.5mm]
 \dfrac{\omega_p^2}{1+\chi} &
 \dfrac{\omega_p\chi\sqrt{2+\chi}}{1+\chi}
 \end{bmatrix}.}
 \label{eq:K-solution}
\end{equation}
\end{theorem}

The result shows explicitly that the two pseudo channels do not generate two
independent scalar corners.  They alter the same second-order closed loop, and
$\chi$ determines how much of the regularization is supplied through the
velocity channel.

If physical acceleration $a$ is injected into the model dynamics, the
acceleration-to-estimated-displacement transfer of the reduced closed loop is
\begin{equation}
 \boxed{
 H_{pa}(s)=
 \frac{1}{1+\chi}
 \frac{1}{s^2+\omega_p\sqrt{2+\chi}\,s+\omega_p^2}.}
 \label{eq:Hpa}
\end{equation}
Define the reconstruction transfer relative to ideal double integration,
\begin{equation}
 G(s):=s^2H_{pa}(s).
 \label{eq:G}
\end{equation}
Then
\begin{equation}
 \lim_{\omega\rightarrow\infty}G(j\omega)=\frac{1}{1+\chi}.
 \label{eq:high-frequency-gain}
\end{equation}
Thus a finite velocity pseudo strength introduces a broadband attenuation
component in addition to the low-frequency drift corner.  This distinction is
consistent with the empirical observation that the position pseudo channel
primarily controls displacement error, whereas the velocity pseudo channel has
stronger interaction with attitude and bias estimation in the complete MEKF.

\section{Physical Displacement MSE}
\label{sec:mse}

The physical objective is written as
\begin{equation}
 J=J_n+J_w,
 \label{eq:J-split}
\end{equation}
where $J_n$ is the displacement variance generated by residual acceleration
noise and $J_w$ is the error caused by reshaping the real wave signal.

\subsection{Residual-noise contribution}

The squared $H_2$ norm of Eq.~\eqref{eq:Hpa} gives the exact reduced noise
variance
\begin{equation}
 \boxed{
 J_n=
 \frac{q}
 {2\omega_p^3(1+\chi)^2\sqrt{2+\chi}}.}
 \label{eq:Jnoise-exact}
\end{equation}
For weak velocity regularization, $\chi\ll1$,
\begin{equation}
 J_n=
 \frac{q}{2\sqrt{2}\,\omega_p^3}
 \left(1-\frac94\chi+O(\chi^2)\right).
 \label{eq:Jnoise-expansion}
\end{equation}
The negative first-order term explains the purpose of the $v=0$ channel in the
reduced model: modest velocity regularization reduces low-frequency stochastic
drift.

\subsection{Wave-distortion contribution}

Let $S_p^{(2)}(\omega)$ be the two-sided physical displacement spectrum and
define
\begin{equation}
 M_0:=\frac{1}{2\pi}\int_{-\infty}^{\infty}
 S_p^{(2)}(\omega)\,d\omega,
 \label{eq:M0}
\end{equation}
\begin{equation}
 M_{-2}:=\frac{1}{2\pi}\int_{-\infty}^{\infty}
 \omega^{-2}S_p^{(2)}(\omega)\,d\omega.
 \label{eq:Mminus2}
\end{equation}
The physical wave error is
\begin{equation}
 J_w=\frac{1}{2\pi}\int_{-\infty}^{\infty}
 |G(j\omega)-1|^2S_p^{(2)}(\omega)\,d\omega.
 \label{eq:Jwave-exact}
\end{equation}
When the regularization corner lies below the principal wave band and
$\chi\ll1$, expansion of Eq.~\eqref{eq:G} gives
\begin{equation}
 \boxed{
 J_w\simeq M_0\chi^2+2M_{-2}\omega_p^2.}
 \label{eq:Jwave-asymptotic}
\end{equation}
The two terms have different physical meanings.  The $\omega_p^2$ term is the
leading displacement distortion caused by the position regularization corner;
the $\chi^2M_0$ term is the leading broadband attenuation caused by the velocity
pseudo channel.

Combining Eqs.~\eqref{eq:Jnoise-expansion} and
\eqref{eq:Jwave-asymptotic} yields the reduced joint objective
\begin{equation}
 \boxed{
 J\simeq
 \frac{q}{2\sqrt{2}\,\omega_p^3}
 \left(1-\frac94\chi\right)
 +2M_{-2}\omega_p^2+M_0\chi^2.}
 \label{eq:J-reduced}
\end{equation}

\section{Joint MSE-Optimal Scaling}
\label{sec:optimal-scaling}

To leading order, optimize Eq.~\eqref{eq:J-reduced} first with respect to
$\omega_p$ and then with respect to $\chi$.

\begin{theorem}[Weak-regularization joint optimum]
In the regime used to obtain Eq.~\eqref{eq:J-reduced}, the stationary point
satisfies
\begin{equation}
 \boxed{
 (\omega_p^\star)^5
 =\frac{3q}{8\sqrt{2}\,M_{-2}},}
 \label{eq:wp-opt}
\end{equation}
and
\begin{equation}
 \boxed{
 \chi^\star
 =\frac{9q}
 {16\sqrt{2}\,M_0(\omega_p^\star)^3}.}
 \label{eq:chi-opt}
\end{equation}
\end{theorem}

For a fixed-shape, period-scaled physical acceleration spectrum, write
\begin{equation}
 S_a^{(2)}(\omega;\sigma_a,\tau)
 =\sigma_a^2\tau\,\Phi_a(\omega\tau),
 \label{eq:Sa-similarity}
\end{equation}
where $\sigma_a$ is a physical acceleration RMS scale.  Since
$S_p^{(2)}=S_a^{(2)}/\omega^4$, similarity gives
\begin{equation}
 M_0\propto\sigma_a^2\tau^4,
 \qquad
 M_{-2}\propto\sigma_a^2\tau^6.
 \label{eq:moment-scaling}
\end{equation}
Substitution into Eqs.~\eqref{eq:wp-opt}--\eqref{eq:chi-opt} gives
\begin{equation}
 \omega_p^\star
 \propto q^{1/5}\sigma_a^{-2/5}\tau^{-6/5},
 \label{eq:wp-scaling}
\end{equation}
\begin{equation}
 \chi^\star
 \propto q^{2/5}\sigma_a^{-4/5}\tau^{-2/5},
 \label{eq:chi-scaling}
\end{equation}
and therefore
\begin{equation}
 \omega_v^\star=\omega_p^\star\sqrt{\chi^\star}
 \propto q^{2/5}\sigma_a^{-4/5}\tau^{-7/5}.
 \label{eq:wv-scaling}
\end{equation}

Using Eq.~\eqref{eq:natural-frequencies}, the corresponding continuous
pseudo-measurement densities are
\begin{equation}
 \rho_p^\star
 \propto q^{1/5}\sigma_a^{8/5}\tau^{24/5},
 \label{eq:rhop-scaling}
\end{equation}
\begin{equation}
 \rho_v^\star
 \propto q^{1/5}\sigma_a^{8/5}\tau^{14/5}.
 \label{eq:rhov-scaling}
\end{equation}
Equation~\eqref{eq:base-density-equivalence} then yields the central base-law
prediction
\begin{equation}
 \boxed{
 r_{p,0}^{\mathrm{base},\star}
 \propto q^{1/10}\sigma_a^{4/5}\tau^{12/5},}
 \label{eq:rp-base-opt}
\end{equation}
\begin{equation}
 \boxed{
 r_{v,0}^{\mathrm{base},\star}
 \propto q^{1/10}\sigma_a^{4/5}\tau^{7/5}.}
 \label{eq:rv-base-opt}
\end{equation}
If the strong direct-acceleration-observation regime makes $q$ approximately
independent of sea amplitude to leading order, the two adaptation laws reduce
to
\begin{equation}
 \boxed{
 r_{p,0}^{\mathrm{base},\star}
 \propto\sigma_a^{4/5}\tau^{12/5},
 \qquad
 r_{v,0}^{\mathrm{base},\star}
 \propto\sigma_a^{4/5}\tau^{7/5}.}
 \label{eq:strong-observation-laws}
\end{equation}
Neither exponent pair is fitted to the implemented schedule.

\section{Comparison with the Implemented Adaptation}
\label{sec:comparison}

Table~\ref{tab:powers} compares the reduced physical-MSE powers with the
implemented OU--II laws.  The per-update rows include the common
$T_S\propto\tau$ cadence normalization of Eq.~\eqref{eq:cadence-normalization}.

\begin{table}[t]
\caption{Implemented and reduced-MSE adaptation powers.}
\label{tab:powers}
\centering
\footnotesize
\begin{tabular}{@{}lll@{}}
\toprule
Quantity & Implemented & Reduced physical MSE \\
\midrule
$r_{p,0}^{\mathrm{base}}$
 & $\sigma_a\tau^2$
 & $\sigma_a^{4/5}\tau^{12/5}$ \\
$r_{v,0}^{\mathrm{base}}$
 & $\sigma_a\tau$
 & $\sigma_a^{4/5}\tau^{7/5}$ \\
$r_p$ per update
 & $\sigma_a\tau^{3/2}$
 & $\sigma_a^{4/5}\tau^{19/10}$ \\
$r_v$ per update
 & $\sigma_a\tau^{1/2}$
 & $\sigma_a^{4/5}\tau^{9/10}$ \\
\bottomrule
\end{tabular}
\end{table}

The discrepancy is the same in both channels.  After absorbing fixed
coefficients,
\begin{equation}
 \boxed{
 \frac{r_{p,\mathrm{MSE}}}{r_{p,\mathrm{impl}}}
 \propto
 \frac{r_{v,\mathrm{MSE}}}{r_{v,\mathrm{impl}}}
 \propto
 q^{1/10}\sigma_a^{-1/5}\tau^{2/5}.}
 \label{eq:common-correction}
\end{equation}
For approximately constant $q$,
\begin{equation}
 \frac{r_{\mathrm{MSE}}}{r_{\mathrm{impl}}}
 \propto\left(\frac{\tau^2}{\sigma_a}\right)^{1/5}.
 \label{eq:fifth-root-correction}
\end{equation}
The fifth root makes the shape difference weak over a finite operating range.

The common wave-band tuner has the following OU--II fixed operating points for
the four JONSWAP calibration records:
$(\tau,\sigma_{aw})=(1.247,0.333)$,
$(2.179,0.683)$,
$(3.590,1.061)$, and $(4.222,1.340)$ in SI units.  Because the mapping from the
physical acceleration RMS to $\sigma_{aw}$ is a fixed multiplicative factor, it
cancels in a normalized comparison.  Normalizing
Eq.~\eqref{eq:fifth-root-correction} at the $H_s=\SI{1.5}{m}$ point gives
Table~\ref{tab:envelope}.

\begin{table}[t]
\caption{Reduced-MSE/common implemented shape ratio over the calibrated
JONSWAP envelope, normalized at $H_s=\SI{1.5}{m}$.}
\label{tab:envelope}
\centering
\footnotesize
\begin{tabular}{@{}lrr@{}}
\toprule
$H_s$ (m) & MSE/implemented & Difference \\
\midrule
0.27 & 0.924 & $-7.6\%$ \\
1.50 & 1.000 & $0.0\%$ \\
4.00 & 1.118 & $+11.8\%$ \\
8.50 & 1.139 & $+13.9\%$ \\
\bottomrule
\end{tabular}
\end{table}

Thus the first-principles power laws produce pseudo standard deviations within
approximately $-8\%$ to $+14\%$ of the implemented shape over the calibrated
JONSWAP range when both are given the same nominal anchor.  This comparison is
in configured pseudo-measurement standard deviation; it does not imply the same
percentage change in displacement RMS.

\section{A Stronger Prediction: the Position-to-Velocity Ratio}
\label{sec:ratio}

The relative channel scaling is more tightly identified than the common
amplitude scale because the unknown $q$ cancels.

\begin{corollary}[Optimal pseudo-channel ratio]
At the weak-regularization optimum,
\begin{equation}
 \boxed{
 \left(\frac{r_{p,0}^{\star}}{r_{v,0}^{\star}}\right)^2
 =\frac{3}{2}\frac{M_{-2}}{M_0}.}
 \label{eq:ratio-exact}
\end{equation}
\end{corollary}

For a fixed normalized sea shape,
$M_{-2}/M_0\propto\tau^2$, hence
\begin{equation}
 \boxed{
 \frac{r_{p,0}^{\star}}{r_{v,0}^{\star}}\propto\tau.}
 \label{eq:ratio-tau}
\end{equation}
This is exactly the relative period scaling of the implemented pair,
\begin{equation}
 \frac{r_{p,0}^{\mathrm{impl}}}{r_{v,0}^{\mathrm{impl}}}
 =\frac{c_p}{c_v}\tau
 =0.500\,\tau.
 \label{eq:implemented-ratio}
\end{equation}

For an ideal infinite-band Pierson--Moskowitz spectrum, the spectral moments can
be integrated analytically and Eq.~\eqref{eq:ratio-exact} gives
\begin{equation}
 \frac{r_{p,0}^{\star}}{r_{v,0}^{\star}}
 =\sqrt{\frac{3}{4\pi}}\,\tau
 \simeq0.489\,\tau,
 \label{eq:pm-ratio}
\end{equation}
when $\tau=T_z/2$.  For an ideal JONSWAP spectrum with peak-enhancement factor
$\gamma=3.3$, numerical moment integration gives approximately
\begin{equation}
 \frac{r_{p,0}^{\star}}{r_{v,0}^{\star}}
 \simeq0.465\,\tau.
 \label{eq:jonswap-ratio}
\end{equation}
The actual simulator uses finite $0.02$--$0.8$ Hz spectra.  Evaluating the same
moment ratio on those underlying first-order spectra over the four reference
periods gives coefficient ratios approximately $0.43$--$0.49$ across the
JONSWAP and PM families.  The implemented value $0.500$ is therefore slightly
looser on the position channel relative to velocity, but it is close to the
reduced physical-MSE prediction without having been obtained from
Eq.~\eqref{eq:ratio-exact}.

This agreement is more informative than agreement of either absolute
coefficient alone.  The ratio follows from the spectral distribution of the
physical displacement and is independent of the assumed residual-acceleration
intensity $q$ and of the common cadence normalization.

\section{Interpretation for OU--II Adaptation}
\label{sec:interpretation}

The reduced theory suggests that the empirical OU--II law has captured the
geometry of the two-channel regularizer more accurately than its simple powers
might imply.  Three observations are notable.

First, the physical-MSE amplitude exponent is $4/5$, already close to the
implemented linear dependence on $\sigma_{aw}$.  As in the OU--III analysis,
a modest difference in exponent produces only a weak schedule difference over
the tested sea-state range.

Second, the theoretical period powers exceed the implemented powers by the
same $2/5$ in both base channels:
\begin{equation}
 \frac{12}{5}-2=\frac{7}{5}-1=\frac25.
\end{equation}
Consequently, the theory does not ask for a different relative period law
between the two pseudo channels.  It asks only for a common multiplicative
shape correction.

Third, the exact relative law
$r_{p,0}/r_{v,0}\propto\tau$ is already present in the implementation.  The
current coefficient ratio $0.65/1.30=0.5$ is also near the spectral value
predicted by the reduced MSE.  This explains why independent empirical sweeps
identify both coefficients sharply while still finding their best region close
to the existing pair.

\section{Testable Predictions and Next Study}
\label{sec:next-study}

The reduced analysis produces a direct experiment rather than another free
exponent search.  Compare the existing laws
\begin{equation}
 r_{p,0}\propto\sigma_{aw}\tau^2,
 \qquad
 r_{v,0}\propto\sigma_{aw}\tau
 \label{eq:current-test-law}
\end{equation}
against
\begin{equation}
 \boxed{
 r_{p,0}\propto\sigma_{aw}^{4/5}\tau^{12/5},
 \qquad
 r_{v,0}\propto\sigma_{aw}^{4/5}\tau^{7/5}.}
 \label{eq:theory-test-law}
\end{equation}
Both candidate pairs should use the same nominal calibration record and the
same one-parameter common gain normalization.  The ratio between the two
channels should not be re-fitted independently unless the experiment is
explicitly testing the spectral-ratio prediction.  The primary comparison
should use paired multi-seed stationary seas, with the controlled transition
reported separately because the existing tuning studies show a distinct
stationary--transient tradeoff.

A second test can remove the empirical channel ratio entirely by using
Eq.~\eqref{eq:ratio-exact} with the measured or versioned displacement spectrum
to set $r_{p,0}/r_{v,0}$ and fitting only one overall regularization scale.  If
that constrained schedule performs comparably to the present two-coefficient
law, it would reduce the steady-state OU--II regularizer from two fitted gains
to one global normalization plus measured sea moments.

The next analytical layer should then use the complete OU--II MEKF
linearization.  This is especially important for the velocity pseudo channel:
repository coefficient sweeps show that $r_{v,0}$ affects roll/pitch more
strongly than the displacement-only reduction predicts.  A full-state
physical-$H_2$ calculation should retain attitude/gravity leakage, gyroscope
and accelerometer bias, latent OU acceleration, and both pseudo channels before
claiming that Eq.~\eqref{eq:theory-test-law} is the final optimal law.

\section{Limitations}
\label{sec:limitations}

The derivation is a low-frequency reduced model.  It assumes that the residual
acceleration error entering the integration chain can be approximated as white
over the regularization band and that the pseudo-measurement corner is below
the dominant physical wave band.  The power laws are obtained from the
$\chi\ll1$ and weak-distortion expansion; the exact transfer in
Eq.~\eqref{eq:Hpa} remains valid outside that asymptotic step, but the closed
form MSE optimum does not.

The displacement objective omits attitude, magnetic, gyroscope-bias, and
accelerometer-bias penalties.  In the full MEKF, both pseudo-measurements act
through covariance cross-couplings, and the velocity channel in particular can
change attitude estimation.  The use of a scalar $q$ also suppresses the large
axis dependence caused by gravity leakage.  Finally, the moment comparison to
JONSWAP and PM uses the underlying first-order finite-band spectra; nonlinear
Stokes harmonics in the simulated trajectories are not represented in those
reference moment constants.

These limitations affect the absolute optimum and may shift the exponents, but
they do not remove the principal structural result that the two pseudo channels
must be optimized jointly and that their leading relative period dependence is
set by physical spectral moments.

\section{Conclusion}

A joint physical-MSE analysis of the OU--II displacement and velocity zero
pseudo-measurements gives a closed-form reduced Riccati model and a direct
adaptation prediction.  In the strong direct-acceleration-observation limit,
the base pseudo standard deviations scale as
\begin{equation}
 \boxed{
 r_{p,0}^{\star}\propto\sigma_a^{4/5}\tau^{12/5},
 \qquad
 r_{v,0}^{\star}\propto\sigma_a^{4/5}\tau^{7/5}.}
\end{equation}
The implemented $\sigma_a\tau^2$ and $\sigma_a\tau$ laws differ from these by
the same fifth-root correction and remain within roughly $-8\%$ to $+14\%$ of
the theoretical shape over the calibrated JONSWAP operating range after a
common nominal normalization.

More strongly, the reduced optimum predicts
\begin{equation}
 \boxed{
 \left(r_{p,0}^{\star}/r_{v,0}^{\star}\right)^2
 =(3/2)M_{-2}/M_0,}
\end{equation}
which gives $r_{p,0}^{\star}/r_{v,0}^{\star}\propto\tau$.  The implemented
coefficient ratio produces $0.500\tau$, close to the approximately
$0.43$--$0.49\tau$ range obtained from the finite-band reference spectra.
Thus the existing two-channel law lies near a first-principles physical-MSE
schedule in both exponent space and relative channel strength.  The appropriate
next step is a paired multi-seed comparison of the analytical and implemented
power laws followed by a full-state physical-$H_2$ analysis of the residual
attitude and bias coupling.

\begin{thebibliography}{9}
\bibitem{KalmanBucy1961}
R. E. Kalman and R. S. Bucy, ``New results in linear filtering and prediction
theory,'' \emph{Journal of Basic Engineering}, vol. 83, no. 1, pp. 95--108,
1961.

\bibitem{Jazwinski1970}
A. H. Jazwinski, \emph{Stochastic Processes and Filtering Theory}.
New York, NY, USA: Academic Press, 1970.
\end{thebibliography}

\end{document}
